% element: 10-s019:e04
\BeleskeID{10-s019}
Jednačine struja koje važe pre izbora režima jesu
\[I_{B1}=\frac{V_I-V_{BE1}}{R_B},\qquad I_{C1}=\frac{V_{CC}-V_{CE1}}{R_C}.\]
U zasićenju koristimo $V_{BES}$ i $V_{CES}$; potrebno je $\beta_FI_{B1}>I_{C1}$. Smanjivanjem ulaznog napona smanjuje se bazna struja. Granica nastupa kada se raspoloživa aktivna kolektorska struja izjednači sa strujom određenom kolektorskim kolom.

Tako dobijeni $V_{I2}$ predstavlja drugi način određivanja $V_{IH}$. Razlika prema $V_{I1}$ potiče samo od zamene $V_{BE}$ sa $V_{BES}$: zajednički član sa otpornicima i napajanjem poništava se pri oduzimanju.
